\begin{table}[!htbp]
\centering\zihao{-4}\setlength{\tabcolsep}{3.5pt}
\caption{附件4全部20条样本的分类贡献与输入来源；区间为从零开始的对齐特征槽位$[a,b)$}
\label{tab:q3_attachment4_explanations}
\begin{tabular}{@{}lrrrlcl@{}}
\toprule
ID & $\phi_T$ & $\phi_A$ & $\phi_V$ & 主导模态 & 关键区间 & 输入来源\\
\midrule
\texttt{\detokenize{01}} & +0.4448 & +0.0283 & +0.0097 & 文本 & $[17,24)$ & 原文字符级\\
\texttt{\detokenize{02}} & -0.1971 & +0.2277 & +0.8294 & 视觉 & $[1,7)$ & 未对齐特征行级\\
\texttt{\detokenize{03}} & +0.6919 & -0.1267 & -0.0492 & 文本 & $[0,10)$ & 原文字符级\\
\texttt{\detokenize{04}} & +1.9249 & +0.1198 & -0.0096 & 文本 & $[20,28)$ & 原文字符级\\
\texttt{\detokenize{05}} & +1.1092 & -0.3336 & +0.6619 & 文本 & $[0,5)$ & 原文字符级\\
\texttt{\detokenize{06}} & +2.2790 & -0.2529 & -0.5109 & 文本 & $[14,25)$ & 原文字符级\\
\texttt{\detokenize{07}} & +1.6925 & -0.3365 & +0.3790 & 文本 & $[13,28)$ & 原文字符级\\
\texttt{\detokenize{08}} & +2.4908 & -0.2620 & -0.5967 & 文本 & $[8,12)$ & 原文字符级\\
\texttt{\detokenize{09}} & +4.0872 & +0.4352 & +0.1656 & 文本 & $[2,8)$ & 原文字符级\\
\texttt{\detokenize{10}} & +3.6626 & +0.4324 & +0.0621 & 文本 & $[21,31)$ & 原文字符级\\
\texttt{\detokenize{11}} & +2.0970 & +0.2465 & +0.7610 & 文本 & $[1,6)$ & 原文字符级\\
\texttt{\detokenize{12}} & +2.1874 & -0.0032 & +0.8495 & 文本 & $[31,44)$ & 原文字符级\\
\texttt{\detokenize{13}} & +0.5928 & +0.0180 & +0.0000 & 文本 & $[0,10)$ & 原文字符级\\
\texttt{\detokenize{14}} & +1.5361 & +0.0546 & -1.0233 & 文本 & $[1,12)$ & 原文字符级\\
\texttt{\detokenize{15}} & +2.9847 & -0.3149 & -0.0090 & 文本 & $[14,21)$ & 原文字符级\\
\texttt{\detokenize{16}} & +3.8314 & +0.4557 & +0.3986 & 文本 & $[8,12)$ & 原文字符级\\
\texttt{\detokenize{17}} & +3.4202 & -0.1421 & +0.1781 & 文本 & $[9,17)$ & 原文字符级\\
\texttt{\detokenize{18}} & +0.8431 & -0.0933 & -0.0955 & 文本 & $[6,21)$ & 原文字符级\\
\texttt{\detokenize{19}} & +0.6310 & -0.4071 & -0.1692 & 文本 & $[29,41)$ & 原文字符级\\
\texttt{\detokenize{20}} & +2.0873 & -0.2851 & +0.1195 & 文本 & $[0,14)$ & 原文字符级\\
\bottomrule
\end{tabular}
\par\vspace{10bp}\begin{minipage}{\textwidth}\songti\normalsize\setlength{\parindent}{2em}
\QThreeParagraph 两个案例的回溯终点不同：样本14可以直接读取原文字符，样本02当前可查到未对齐视觉特征行。同一遮挡方法给出特征区间，而能够进一步阅读或展示何种输入，取决于该模态保存的来源记录。

\QThreeParagraph 本问固定问题二的预测模型，以模态组合输出计算三模态贡献与两两交互，再用连续窗口遮挡定位局部影响。随机等长窗口、删除曲线和跨种子比较分别考察局部响应、累计影响及初始化变化，使解释结果可以结合模型的实际输出评估。

\QThreeParagraph 附件4的20条样本均得到类别、连续强度、模态贡献和关键区间，并按已有来源记录连接到原文或未对齐特征行。逐样本的同向与反向贡献、不同主导模态及不同位置，共同呈现固定模型使用三路输入的差异，也为检查具体预测提供了可复查的记录。

\end{minipage}
\end{table}
